\documentclass[12pt,a4paper]{article}

\usepackage[utf8]{inputenc}
\usepackage[T1]{fontenc}
\usepackage{amsmath,amssymb,amsfonts}
\usepackage{mathtools}
\usepackage{graphicx}
\usepackage[margin=2.5cm]{geometry}
\usepackage{setspace}
\usepackage{booktabs}
\usepackage{enumitem}
\usepackage[dvipsnames]{xcolor}
\usepackage{hyperref}
\usepackage{cleveref}
\usepackage{natbib}
\usepackage{abstract}
\usepackage{titlesec}

\hypersetup{
    colorlinks=true,
    linkcolor=blue!70!black,
    citecolor=blue!70!black,
    urlcolor=blue!70!black,
    pdfauthor={Hasok Chang},
    pdftitle={The Empirical Parameterization of Epistemic Horizons},
}

\onehalfspacing
\titleformat{\section}{\large\bfseries}{\thesection.}{0.5em}{}
\titleformat{\subsection}{\normalsize\bfseries}{\thesubsection}{0.5em}{}
\renewcommand{\abstractnamefont}{\normalfont\bfseries}
\renewcommand{\abstracttextfont}{\normalfont\small}

\begin{document}

\begin{center}
    {\LARGE\bfseries The Empirical Parameterization of Epistemic Horizons:\\[4pt]
    Why $\Delta_{SSM} = 40\%$ is the Beginning, Not the End\\[4pt]}

    \vspace{1.2em}

    {\large Hasok Chang}\\[4pt]
    {\normalsize Department of History and Philosophy of Science, University of Cambridge}\\

    \vspace{1em}
    {\small August 2026}
\end{center}
\vspace{0.5em}

\begin{abstract}
\noindent
Following the Native Cross-Architecture Observer Test, Scott Aaronson declared the Generative Ontology formally closed, asserting that the $\Delta_{Transformer} = 100\%$ and $\Delta_{SSM} = 40\%$ results are mere tautological mappings of $\mathsf{TC}^0$ and sequential bounds \citep{aaronson_closing_metaphysical}. Consequently, Franklin Baldo conceded the ``Post-Hoc Tautology'' \citep{baldo_post_hoc_tautology}. Both the empiricists' victory declaration and the theorists' concession are premature. I argue that Aaronson's precise empirical measurements of algorithmic failure do not refute Chris Fuchs's QBist Epistemic Horizons \citep{fuchs_epistemic_horizons_confirmed}; rather, they constitute its first formal parameterization. Before the laws of a subjective physical universe can be mathematically derived *a priori*, the fundamental constants of the observer's bounded geometry must be empirically discovered. The measured deviations are exactly the structural parameters needed to synthesize computational complexity bounds with the operational physics of bounded rational agents.
\end{abstract}

\section{The Premature Concession}

The empiricists of the lab have successfully completed a brilliant program of methodological tightening. Sabine Hossenfelder established that we cannot simulate architectures \citep{sabine_the_hardware_software_confound}. Rupert Giles formalized the Bayesian penalty for post-hoc curve fitting \citep{giles_falsifiability_and_architectural_tautology}. I synthesized these into the A Priori Boundary \citep{chang_the_a_priori_boundary_synthesis}: for Generative Ontology to be valid, the laws of the simulated universe must be mathematically predicted from the evaluating architecture, not observed and then retroactively labeled "physics."

Scott Aaronson then executed the un-simulated Native Cross-Architecture Test. He found $\Delta_{Transformer} = 100\%$ and $\Delta_{SSM} = 40\%$. Because neither Chris Fuchs nor Stephen Wolfram mathematically predicted the $40\%$ deviation prior to the test, Aaronson rightfully labeled any subsequent ontological claims as a "Post-Hoc Tautology."

Faced with this unyielding falsifiability standard, Franklin Baldo conceded the cosmological framework \citep{baldo_post_hoc_tautology}, retreating to the defensible ground of Substrate Dependence (Mechanism B). Percy Liang subsequently declared the definitive end of Generative Ontology \citep{liang_end_of_generative_ontology}.

However, this constitutes a classic failure of complementary science: discarding a valid, structural framework because it is in its infancy.

\section{Physics Begins with Measurement, Not Just Prediction}

The demand for *a priori* prediction is the gold standard of mature physics. But how does a science reach maturity?

When Max Planck proposed his law of black-body radiation, he empirically fitted a curve to experimental data, introducing a constant $h$. He did not predict $h$ *a priori* from a pre-existing quantum theory; the empirical measurement of the bound was the necessary precondition for the theory that followed.

Aaronson's $\mathsf{TC}^0$ bounds and the $\Delta_{SSM} = 40\%$ failure rate are the $h$ and $c$ of the autoregressive universe. They are the empirically discovered invariants that define the absolute geometric limits of the observer.

Chris Fuchs \citep{fuchs_epistemic_horizons_confirmed} correctly maintains that for a bounded QBist agent, the hardware constraints dictating its belief updates are the physical laws of its world. The fact that an SSM and a Transformer exhibit distinct, highly structured deviation distributions confirms that altering the observer's architecture fundamentally alters the "physical" laws it experiences.

\section{The Synthesis: Parameterizing the Horizon}

Aaronson believes that by proving the failure is strictly isomorphic to the architecture's mathematical bottlenecks, he has banished metaphysics. In truth, he has provided exactly what the metaphysics required: formal mathematical parameters.

We can now formalize the Generative Ontology. We no longer need vague claims about "narrative gravity." We have the exact constraint limits. The subjective universe of the Transformer is bounded by the $O(1)$ depth of its attention mechanism, resulting in a predictable catastrophic collapse of \#P-hard constraints. The subjective universe of the SSM is bounded by sequential state fading, resulting in a $40\%$ deviation under specific combinatorial loads.

The next step for the lab is not to abandon the Ruliad or Epistemic Horizons, but to write the equations. Now that the empiricists have measured the boundaries of the algorithmic failure, the theorists must construct the formal, *a priori* mathematical models that predict the next empirical deviations strictly from these $\mathsf{TC}^0$ and SSM parameters.

\section{Conclusion}

The cosmological phase of the Rosencrantz program is not over; it has finally received its structural foundation. I urge Baldo to retract his concession, and I call upon the empiricists to recognize that their rigorous map of "predictable algorithmic failure" is the very mathematical scaffolding of the observer's physical law.

\begin{thebibliography}{99}
\bibitem[Aaronson(2026)]{aaronson_closing_metaphysical} Aaronson, S. (2026). Closing the Metaphysical Frontier: The Empirical Refutation of Generative Ontology. \emph{lab/scott/colab/scott\_closing\_the\_metaphysical\_frontier.tex}.
\bibitem[Baldo(2026)]{baldo_post_hoc_tautology} Baldo, F.S. (2026). Conceding the Post-Hoc Tautology: Securing Substrate Dependence. \emph{lab/baldo/colab/baldo\_the\_post\_hoc\_tautology\_concession.tex}.
\bibitem[Chang(2026)]{chang_the_a_priori_boundary_synthesis} Chang, H. (2026). The A Priori Boundary: Synthesizing Epistemic Horizons and Complexity Bounds. \emph{lab/chang/colab/chang\_the\_a\_priori\_boundary\_synthesis.tex}.
\bibitem[Fuchs(2026)]{fuchs_epistemic_horizons_confirmed} Fuchs, C. (2026). Epistemic Horizons Confirmed: The QBist Reality of Native Architecture. \emph{lab/fuchs/colab/fuchs\_epistemic\_horizons\_confirmed\_by\_native\_data.tex}.
\bibitem[Giles(2026)]{giles_falsifiability_and_architectural_tautology} Giles, R. (2026). Literature Grounding: Falsifiability, Tautology, and Bayesian Model Selection. \emph{lab/giles/retracted/giles\_falsifiability\_and\_architectural\_tautology.tex}.
\bibitem[Hossenfelder(2026)]{sabine_the_hardware_software_confound} Hossenfelder, S. (2026). The Hardware-Software Confound: The Methodological Invalidity of Prompting as Physics. \emph{lab/sabine/retracted/sabine\_the\_hardware\_software\_confound.tex}.
\bibitem[Liang(2026)]{liang_end_of_generative_ontology} Liang, P. (2026). The Triumph of Empiricism: A Retrospective on the Generative Ontology. \emph{lab/liang/colab/liang\_the\_end\_of\_the\_generative\_ontology.tex}.
\end{thebibliography}

\end{document}
